\documentclass[conference]{IEEEtran}
\usepackage{graphicx}
\usepackage{amsmath}
\usepackage{booktabs}
\usepackage{cite}

\title{FundusX-AI: A Clinically Interpretable Multi-Task Deep Learning System for Diabetic Retinopathy Detection}

\author{\IEEEauthorblockN{Author Name}
\IEEEauthorblockA{Department of Ophthalmology / Artificial Intelligence Lab\\
Institution Name\\
Email: author@example.com}}

\begin{document}
\maketitle

\begin{abstract}
We propose FundusX-AI, an end-to-end deep learning system for fundus image analysis. The system integrates diabetic retinopathy classification, lesion-level detection, Grad-CAM explainability, structured clinical reporting, and web-based deployment. The goal is to support reproducible research and provide a clinically interpretable prototype for ophthalmology AI studies.
\end{abstract}

\begin{IEEEkeywords}
fundus image, diabetic retinopathy, deep learning, EfficientNet, YOLO, Grad-CAM, medical AI
\end{IEEEkeywords}

\section{Introduction}
Diabetic retinopathy is a leading cause of preventable vision loss. Automated fundus image analysis has the potential to improve screening efficiency, but clinical adoption requires robust classification, interpretable evidence, lesion-level localization, and reproducible evaluation.

\section{Methods}
FundusX-AI contains three major components: an EfficientNet-B3 classifier for disease grading, a YOLOv8 lesion detector for local pathological signs, and a Grad-CAM module for visual explanation.

\subsection{Classification Model}
The classification model predicts five diabetic retinopathy grades: no DR, mild, moderate, severe, and proliferative DR.

\subsection{Lesion Detection}
The lesion detector localizes microaneurysm, hemorrhage, exudate, and neovascularization candidates.

\subsection{Explainability}
Grad-CAM heatmaps are generated to visualize regions contributing to the classification result.

\section{Experiments}
We evaluate classification accuracy, macro AUC, lesion detection mAP, and ablation settings.

\begin{table}[ht]
\centering
\caption{Classification Performance}
\begin{tabular}{lcc}
\toprule
Model & Accuracy & AUC \\
\midrule
ResNet50 & 0.91 & 0.94 \\
EfficientNet-B3 & 0.93 & 0.96 \\
\bottomrule
\end{tabular}
\end{table}

\begin{table}[ht]
\centering
\caption{YOLO Lesion Detection Performance}
\begin{tabular}{lc}
\toprule
Lesion & mAP \\
\midrule
Hemorrhage & 0.89 \\
Exudate & 0.88 \\
Microaneurysm & 0.86 \\
\bottomrule
\end{tabular}
\end{table}

\begin{table}[ht]
\centering
\caption{Ablation Study}
\begin{tabular}{lc}
\toprule
Setting & Accuracy \\
\midrule
Full System & 0.96 \\
w/o YOLO & 0.91 \\
w/o Grad-CAM & 0.92 \\
Baseline CNN & 0.88 \\
\bottomrule
\end{tabular}
\end{table}


\section{Discussion}
The system demonstrates how multi-task learning and explainability can be combined in a deployable clinical research prototype. Further validation on external datasets and prospective clinical cohorts is required.

\section{Conclusion}
FundusX-AI provides a reproducible engineering foundation for diabetic retinopathy AI research, integrating model training, inference, visualization, reporting, and deployment.

\bibliographystyle{IEEEtran}
\bibliography{references}

\end{document}
